% Generated draft: requires mathematical and source review.
\documentclass[11pt]{article}
\usepackage[margin=25mm]{geometry}
\usepackage{fontspec}
\setmainfont{Latin Modern Roman}
\usepackage{amsmath,amssymb,mathtools,unicode-math}
\setmathfont{Latin Modern Math}
\usepackage{enumitem,needspace,xurl,hyperref}
\setlength{\parindent}{0pt}
\setlength{\parskip}{5pt}
\setlength{\emergencystretch}{3em}
\newcommand{\N}{\mathbb{N}}
\newcommand{\Z}{\mathbb{Z}}
\newcommand{\Q}{\mathbb{Q}}
\newcommand{\R}{\mathbb{R}}
\newcommand{\C}{\mathbb{C}}
\newcommand{\F}{\mathbb{F}}
\newcommand{\A}{\mathbb{A}}
\newcommand{\PP}{\mathbb{P}}
\newcommand{\E}{\mathbb{E}}
\newcommand{\eps}{\varepsilon}
\newcommand{\dd}{\,\mathrm{d}}
\begin{document}
\section*{UnsolvedMath ID: 3800002}
\textbf{Source problem code:} catalogue entry
\subsection{Definitions and mathematical statement}
Let \(Q=[0,1]^3\subseteq\R^3\). A tetrahedral triangulation of \(Q\) is a finite collection \(\mathcal T\) of nondegenerate straight Euclidean tetrahedra contained in \(Q\), whose union is \(Q\), such that any two distinct members intersect in a common face, edge, vertex, or the empty set. Vertices may be added on the boundary and in the interior; no restriction to the eight cube vertices is stated. For an edge \(e\) of \(T\in\mathcal T\), let \(\theta(T,e)\in(0,\pi)\) be the interior dihedral angle between the two incident triangular faces. The question is
\[
\exists\mathcal T\text{ triangulating }Q\quad
\forall T\in\mathcal T\ \forall e\text{ an edge of }T,\qquad
\theta(T,e)<\frac\pi2.
\]
The inequalities are strict; a nonobtuse triangulation alone does not answer the question. The supplied research summary reports an affirmative solution by two constructions in the references below. It also reports nonexistence of acute triangulations of the \(n\)-cube for \(n\ge4\), as an ancillary higher-dimensional result rather than an additional hypothesis on the three-dimensional question. These are supplied literature assessments, not independently verified here.
\subsection{Short English statement}
Can the cube be divided into tetrahedra meeting face to face, with every interior dihedral angle strictly less than \(90\) degrees? The supplied registry records an affirmative literature solution.
\subsection{Sources}
\begin{itemize}
\item Local UnsolvedMath registry, permanent ID \(3800002\), exact source code \texttt{catalogue entry}, ``Acute triangulation of the cube.'' Proposer field: Jeff Erickson and cited proposers.
\item Erickson, \emph{Open Problems in Computational Geometry}, item ``Acute triangulation of the cube''; \url{https://jeffe.cs.illinois.edu/open/}. HTML extraction, recorded access 2026-07-29.
\item E. Kopczyński, I. Pak, P. Przytycki, ``Acute triangulations of polyhedra and \(\R^n\),'' \emph{Combinatorica} \textbf{32} (2012), 583--608; arXiv:1110.6084, as supplied in the research summary. The summary attributes cube and regular-octahedron constructions and the \(n\ge4\) cube obstruction to this work.
\item E. VanderZee, A. N. Hirani, D. Guoy, E. A. Ramos, ``A dihedral acute triangulation of the cube,'' \emph{Computational Geometry} \textbf{46} (2013), 492--506; DOI 10.1016/j.comgeo.2010.09.002. The supplied summary reports seeing the abstract's existence statement.
\end{itemize}
\end{document}
